\chapter{Wymogi formalne}\label{ch:wymogi}

Zestawienie wymagań obowiązujących pracę licencjacką na kierunku elektroniczne przetwarzanie informacji, z podaniem źródła każdego z nich. Dwa dokumenty Instytutu uzupełniają się i miejscami rozchodzą -- przy rozbieżności o \wyroznienie{układzie pracy} rozstrzygają Standardy \parencite{standardyEPI}, a o \wyroznienie{zapisie bibliografii} Instrukcja \parencite{instrukcjaISI}. W tabelach skrót S oznacza Standardy, a skrót I -- Instrukcję.

\section{Skład}

\begin{tabelaepi}{Wymogi składu}{tab:sklad}
\begin{tabular}{@{}p{8.4cm}cp{3.4cm}@{}}
\toprule
Wymóg & Źródło & Kto realizuje \\
\midrule
Pismo 12 pkt w tekście głównym & S & klasa \\
Interlinia 1,5 & S & klasa \\
Marginesy 2,5 cm z każdej strony & S, I & klasa \\
Przy oprawie margines wewnętrzny szerszy kosztem zewnętrznego, nie mniej niż 1,5 cm & I & opcja klasy \\
Tytuły rozdziałów tym samym krojem, większe i wytłuszczone & S & klasa \\
Tytuły rozdziałów 18 pkt, wyśrodkowane & I & klasa \\
Tytuły podrozdziałów 14 pkt, wytłuszczone, do lewej & I & klasa \\
Numeracja rozdziałów wielopoziomowa & S, I & klasa \\
Wstęp i Zakończenie bez numeracji, ale w spisie treści & I & klasa \\
Każdy element ze spisu treści od nowej strony & S, I & klasa \\
Wszystkie strony numerowane & S & klasa \\
Numer strony w prawym dolnym rogu & I & klasa \\
Strona tytułowa liczona, bez numeru & S, I & klasa \\
Tekst wyjustowany & S, I & klasa \\
\bottomrule
\end{tabular}
\end{tabelaepi}

Wszystkie pozycje realizuje klasa dokumentu bez Twojego udziału. Nie zmieniaj ustawień składu w preambule -- rozbieżność z wymogami obciąża pracę, a nie klasę.

\section{Objętość}

Tekst od pierwszej strony Wprowadzenia do ostatniej strony Podsumowania nie może przekraczać \wyroznienie{czterdziestu stron znormalizowanych, czyli 72 000 znaków ze spacjami} \parencite{instrukcjaISI}. Do limitu nie wliczają się: strona tytułowa, strony informacyjne, spis treści, wykaz źródeł, bibliografia, spisy, indeks i aneksy.

Kontrolę wykonuje skrypt \plik{tools/policz\_znaki.R}, uruchamiany z katalogu wzoru. Liczy znaki po usunięciu poleceń, listingów, wzorów i tabel, więc daje oszacowanie od góry.

\section{Struktura}

Kolejność elementów we wzorze łączy wymagania obu dokumentów: strona tytułowa, opcjonalna strona informacyjna z abstraktem i słowami kluczowymi, spis treści, Wprowadzenie, rozdziały numerowane, Podsumowanie, wykaz źródeł, bibliografia, spis ilustracji, indeks nazwisk, aneksy.

Standardy opisują treść rozdziałów zasadniczych jako logikę aplikacji -- schemat ogólny, moduły, opis interfejsu ze zrzutami ekranu -- oraz implementację, czyli format danych, strukturę bazy, sposób realizacji modułów, biblioteki zewnętrzne i narzędzia. Wzór realizuje to w trzech rozdziałach.

\subsection{Strona informacyjna}

Wymóg Instrukcji, włączany opcją klasy. Zawiera opis bibliograficzny pracy w postaci obejmującej nazwisko i imię autora, rok obrony, tytuł, oznaczenie rodzaju pracy, promotora, miejsce, jednostkę i liczbę stron. Dalej abstrakt liczący \wyroznienie{od 1000 do 1500 znaków ze spacjami}, przedstawiający cel, przedmiot, metody i kluczowe wyniki, napisany językiem skondensowanym i bez ogólników. Na końcu \wyroznienie{około pięciu słów kluczowych} w porządku alfabetycznym, rozdzielonych przecinkami.

Strona angielska powstaje automatycznie po wypełnieniu angielskich metadanych i zawiera dokładnie te same informacje w tym samym układzie.

\section{Cytowanie}

Standardy wymieniają pięć sytuacji, w których wskazanie źródła jest obowiązkowe:

\begin{enumerate}
\item cytat bezpośredni, czyli dosłownie przytoczone cudze słowa;
\item graficzna forma cytatu -- cudze tabele, rysunki, zestawienia;
\item cytat znany pośrednio, z tekstu jeszcze innego autora;
\item zebrane przez kogoś informacje: ankiety, dane statystyczne, dokumenty, materiały ze stron internetowych;
\item cudza teza, argument, opinia, idea, interpretacja, unikalne ujęcie tematu.
\end{enumerate}

Każde przytoczenie umieszczasz w cudzysłowie polskim, a cytat w cytacie w cudzysłowie wewnętrznym. Wewnątrz cytatu obowiązuje pisownia cytowanego autora; zmianę odnotowujesz w nawiasie kwadratowym po cudzysłowie zamykającym. Tłumaczenie własne oznaczasz w nawiasie kwadratowym, a cytat w oryginale przywołujesz w przypisie dolnym. Fragment powyżej dwóch lub trzech zdań wyodrębniasz graficznie.

Przytoczenie musi być \wyroznienie{wkomponowane w tekst}: wprowadzone lub domknięte komentarzem, który sygnalizuje obecność cudzego słowa. Samo wklejenie z przypisem nie wystarcza. Przy powoływaniu się na cudzą opinię obowiązuje, niezależnie od przypisu, wprowadzenie z informacją, do kogo dany pogląd należy.

Brak zasygnalizowania w tekście, że posługujesz się cudzym słowem, oznacza naruszenie prawa \parencite{prawoAutorskie}. Brak odpowiednio opisanych źródeł może być powodem dyskwalifikacji pracy.

\section{Przypisy i bibliografia}

Obowiązuje system nawiasowy, tak zwany harwardzki: nazwisko, rok, w razie potrzeby strona. Postacie powołań i odpowiadające im polecenia zestawia tabela~\ref{tab:powolania}. Skracanie listy autorów, dodawanie inicjału i sufiksów rocznika wykonuje styl automatycznie -- Twoim zadaniem jest poprawny wpis w pliku bibliograficznym.

Do przypisów dolnych trafiają wyłącznie przypisy dygresyjne, których należy unikać, oraz oryginalne brzmienie tłumaczonych cytatów.

Zasada wiążąca bibliografię z tekstem: \wyroznienie{bibliografia zawiera wyłącznie pozycje faktycznie wykorzystane, do których odsyłają przypisy w tekście}. Pozycje dotyczące tematu, ale niewykorzystane, do bibliografii nie trafiają. Muszą się w niej znaleźć pozycje obcojęzyczne.

\section{Materiał ilustracyjny}

\begin{tabelaepi}{Wymogi wobec materiału ilustracyjnego}{tab:ilustracje}
\begin{tabular}{@{}p{10.6cm}c@{}}
\toprule
Wymóg & Źródło \\
\midrule
Numeracja ciągła przez cały tekst, osobna dla każdego typu & I \\
Tytuł nad ilustracją & I \\
Źródło pod ilustracją & I \\
Oznaczenie opracowania własnego, także na podstawie cudzego materiału & I \\
Pełny opis bibliograficzny dla materiału przejętego & I \\
Interfejs aplikacji zilustrowany zrzutami ekranu & S \\
\bottomrule
\end{tabular}
\end{tabelaepi}

Środowiska wzoru realizują tytuł, numerację i miejsce na źródło automatycznie.

Prawa do ilustracji: bez zgody autora można wykorzystać materiały z domeny publicznej, na licencji Creative Commons lub podobnej oraz zrzuty ekranu z publicznie dostępnych źródeł elektronicznych. Pozostałe wymagają pisemnej zgody właściciela autorskich praw majątkowych.

\section{Zapis w tekście}

\begin{tabelaepi}{Zapis elementów w tekście}{tab:zapiselementow}
\begin{tabular}{@{}p{6.8cm}ll@{}}
\toprule
Element & Zapis & Polecenie \\
\midrule
tytuł czasopisma, książki, zasobu sieciowego & kursywa bez cudzysłowu & \kod{\textbackslash termin} \\
wyraz obcojęzyczny & kursywa & \kod{\textbackslash termin} \\
wyróżnienie treściowe & wytłuszczenie & \kod{\textbackslash wyroznienie} \\
cytat & cudzysłów polski & \kod{\textbackslash enquote} \\
nazwa funkcji & pismo maszynowe & \kod{\textbackslash fun} \\
nazwa argumentu & pismo maszynowe & \kod{\textbackslash argument} \\
nazwa pakietu & pismo maszynowe & \kod{\textbackslash pkg} \\
nazwa pliku & pismo maszynowe & \kod{\textbackslash plik} \\
\bottomrule
\end{tabular}
\end{tabelaepi}

\section{Wymogi projektu dyplomowego}

Poza pracą pisemną Standardy wymagają samodzielnie zaprojektowanej i wykonanej aplikacji z interfejsem internetowym. Zaliczenie seminarium następuje dopiero po przyjęciu projektu \parencite{standardyEPI,prawoSzkolnictwo}.

Aplikacja nastawiona na prezentację informacji musi robić więcej niż wyświetlanie zawartości bazy danych: konieczny jest niestandardowy pomysł prezentacji, dostosowany projekt graficzny i mechanizm interakcji z użytkownikiem.

W serwisie obowiązkowa jest nota z załącznika do Standardów, umieszczona w zakładce o serwisie. Podaje ona autora serwisu, informację, że serwis stanowi integralną część pracy licencjackiej na kierunku elektroniczne przetwarzanie informacji, oraz promotora i jednostkę. Ta sama nota trafia do aneksu pracy.

\section{Kryteria oceny}

Standardy wymieniają wprost, co podlega ocenie promotora i recenzenta: zgodność treści pracy z tematem, strukturę pracy, znajomość stanu badań lub praktyki w zakresie wyznaczonym przez tematykę, poprawność językową i terminologiczną, metodologię projektu, zawartość merytoryczną, poprawność działania oraz estetykę aplikacji, a także dobór i wykorzystanie piśmiennictwa. Recenzje trafiają do Archiwum Prac.

\section{Lista kontrolna przed oddaniem}\label{sec:listakontrolna}

\subsection*{Kompilacja i objętość}

\begin{kontrolna}
\item Pełna kompilacja kończy się bez błędów
\item Brak komunikatów o nierozwiązanych powołaniach i odsyłaczach
\item Objętość mieści się w limicie
\item Listingi nie zawierają znaków spoza ASCII
\end{kontrolna}

\subsection*{Treść}

\begin{kontrolna}
\item Tytuł pracy odpowiada zawartości i zawiera zagadnienie ogólniejsze
\item Cel z Wprowadzenia domknięty w Podsumowaniu
\item Perspektywy rozwoju aplikacji opisane konkretnie
\item Wszystkie przytoczenia wprowadzone komentarzem w tekście
\item Własne ustalenia oddzielone od cudzych na poziomie zdania
\end{kontrolna}

\subsection*{Bibliografia i źródła}

\begin{kontrolna}
\item Każda pozycja bibliografii ma powołanie w tekście
\item Każde powołanie ma pozycję w bibliografii
\item W bibliografii są pozycje obcojęzyczne
\item Wykaz źródeł obejmuje wszystkie użyte pakiety
\item Adresy sieciowe mają datę odczytu
\end{kontrolna}

\subsection*{Elementy końcowe}

\begin{kontrolna}
\item Spis ilustracji zgodny z zawartością pracy
\item Indeks nazwisk obejmuje nazwiska z tekstu i bibliografii
\item Aneksy powiązane z tekstem odsyłaczami
\item Aneks z metadanymi oprogramowania wypełniony
\item Aneks z notą serwisu zgodny z treścią na stronie projektu
\item Oświadczenie o wykorzystaniu sztucznej inteligencji wypełnione, jeżeli dotyczy
\item Wykaz poleceń wydanych narzędziu programistycznemu kompletny
\end{kontrolna}

\subsection*{Projekt}

\begin{kontrolna}
\item Strona projektu działa pod adresem podanym na stronie tytułowej
\item Nota z załącznika obecna w serwisie
\item Repozytorium publiczne i zgodne z adresem ze strony tytułowej
\item Sprawdzenie pakietu przechodzi na trzech systemach operacyjnych
\end{kontrolna}

\subsection*{Redakcja}

\begin{kontrolna}
\item Korekta całości po wydruku, nie na ekranie
\item Tekst wyjustowany, bez wiszących wierszy
\item Konsekwentna terminologia w całej pracy
\item Wyszukanie uwag promotora w plikach źródłowych nie daje wyników
\item Usunięte własne komentarze robocze i wskazówki wzoru
\end{kontrolna}
